%!TEX root = ./main.tex

\section[Origins]{Origins: Resources, Names, Milestones}

% SECTION TIMING: 0:00--9:00 (3 frames).
% COLOR CONVENTION FOR THE WHOLE DECK -- do not break it:
%   ranblue    = the mainstream 2.4/5/6 GHz branch (b/a/g/n/ac/ax/be)
%   coreorange = the 60 GHz ad/ay branch (WiGig)
%   softgreen  = what survives a tradeoff
%   ranpurple  = security / adversary
% LAYOUT RULES for this template:
%   - a frame title must fit ONE line in the title bar: keep it under ~44 characters.
%   - every bullet must fit on one line: ~30 characters in a half-width column.
%     Shorten the text, do not let it wrap.

% Teaching note (180 s): open with the hook -- the Wi-Fi fan icon (drawn above the
% headline, not downloaded) has looked the same since before the students were born,
% yet the radio under it has been rebuilt repeatedly. This table is the organizing
% framework for the entire lecture; the rows reveal one by one (\pause), so name each
% dimension as it appears instead of reading a wall. Ask students to keep a tally as
% the lecture proceeds -- the closing five-questions frame cashes that tally in.
\begin{frame}{Wi-Fi grew by spending new resources}
  \centering
  % The generic Wi-Fi fan glyph, drawn in TikZ (three arcs + a dot), not an image.
  \begin{tikzpicture}[line cap=round]
    \fill[ranblue] (0,0) circle (0.06);
    \foreach \r in {0.2,0.38,0.56}
      \draw[ranblue,line width=1.3pt] (45:\r) arc (45:135:\r);
  \end{tikzpicture}\par
  \vspace{0.05em}
  {\large The icon never changed. The radio under it did --- six ways.\par}

  \vspace{0.1em}
  {\scriptsize\textcolor{softgray}{Preview, not prerequisite --- every row gets its
    own slide today. Keep a tally.}\par}

  \vspace{0.2em}
  % SOURCES for the progressions in this table:
  % https://standards.ieee.org/ieee/802.11/10548/
  % https://standards.ieee.org/ieee/802.11be/7516/
  % https://www.cisco.com/c/en/us/products/collateral/wireless/white-paper-c11-740788.html
  % \visible (not \pause): rows uncover in place, so the layout never jumps.
  {\small
  \begin{tabular}{@{}ll@{}}
    \toprule
    \textbf{Dimension} & \textbf{How Wi-Fi spent it} \\
    \midrule
    Spectrum      & 2.4 $\to$ 5 $\to$ 6 $\to$ \textcolor{coreorange}{60} GHz \\
    \visible<2->{Bandwidth}     & \visible<2->{20 $\to$ 40 $\to$ 80 $\to$ 160 $\to$ 320 MHz} \\
    \visible<3->{Constellation} & \visible<3->{BPSK $\to$ 64 / 256 / 1024 / 4096-QAM} \\
    \visible<4->{Space}         & \visible<4->{SISO $\to$ SU-MIMO $\to$ MU-MIMO} \\
    \visible<5->{Scheduling}    & \visible<5->{one user $\to$ OFDMA resource units} \\
    \visible<6->{Links}         & \visible<6->{one band $\to$ multi-link operation} \\
    \bottomrule
  \end{tabular}\par}

  \vspace{0.2em}
  % 60 GHz (802.11ad, 2012) predates 6 GHz operation (802.11ax-2021 / Wi-Fi 6E).
  {\scriptsize\textcolor{softgray}{Arrows list resources, not chronology:
    \textcolor{coreorange}{60 GHz} (2012) predates 6 GHz (2020s).}\par}

  \vspace{0.2em}
  \takeaway{One family identity; each amendment spends a new dimension.}

  \glossline{QAM = Quadrature Amplitude Modulation \quad MIMO = Multiple Input,
    Multiple Output \quad OFDMA = Orthogonal Frequency-Division Multiple Access}
  \source{https://standards.ieee.org/ieee/802.11/10548/}{IEEE 802.11 standard family; IEEE 802.11be; Cisco 802.11ax white paper}
\end{frame}

% Teaching note (150 s): students constantly mix up "802.11ax" and "Wi-Fi 6". Make the
% naming split explicit once, here, and never again: IEEE names technical amendments,
% the Wi-Fi Alliance sells generation numbers. Point at the orange ad/ay row: it is a
% parallel 60 GHz branch (WiGig), outside the numbered 1-7 sequence. Mention that
% Wi-Fi 1-3 are retroactive informal labels; official Alliance numbering began at Wi-Fi 4.
\begin{frame}{Amendment names vs.\ Wi-Fi generations}
  \centering
  % SOURCE: https://standards.ieee.org/wp-content/uploads/interactive/web/wi-fi-timeline/index.html
  % SOURCE: https://helpfiles.keysight.com/csg/n7637/Content/Main/802.11ad%20Concepts.htm
  \small
  \begin{tabular}{@{}lll@{}}
    \toprule
    \textbf{IEEE amendment} & \textbf{Market name} & \textbf{PHY shorthand} \\
    \midrule
    802.11b  & Wi-Fi 1 & HR-DSSS / CCK \\
    802.11a  & Wi-Fi 2 & OFDM at 5 GHz \\
    802.11g  & Wi-Fi 3 & OFDM at 2.4 GHz \\
    802.11n  & Wi-Fi 4 & HT: MIMO + 40 MHz \\
    802.11ac & Wi-Fi 5 & VHT: wider + 256-QAM \\
    802.11ax & Wi-Fi 6 / 6E & HE: OFDMA + 1024-QAM \\
    802.11be & Wi-Fi 7 & EHT: MLO + 320 MHz \\
    \textcolor{coreorange}{802.11ad / ay} & \textcolor{coreorange}{WiGig} &
      \textcolor{coreorange}{directional 60 GHz} \\
    \bottomrule
  \end{tabular}

  \vspace{0.15em}
  % Defuses the forward-reference wall: column 3 is a promise, not a prerequisite.
  {\scriptsize\textcolor{softgray}{IEEE names amendments; the Wi-Fi Alliance names
    generations (Wi-Fi 1--3 are retroactive, informal). The orange ad/ay row is a
    parallel branch, outside Wi-Fi 1--7. Third column = today's itinerary ---
    every acronym gets its own slide. Read columns 1--2 now.}\par}

  \vspace{0.15em}
  \takeaway{One radio, two naming systems: engineering and marketing.}

  \glossline{HR = High Rate \quad HT/VHT/HE/EHT = High / Very High / High-Efficiency
    / Extremely High Throughput \quad MLO = Multi-Link Operation}
  \source{https://standards.ieee.org/wp-content/uploads/interactive/web/wi-fi-timeline/index.html}{IEEE 802.11 interactive timeline; Keysight 802.11ad concepts}
\end{frame}

% Teaching note (180 s): problem before solution -- before 1999 there was no consumer
% Wi-Fi, and each milestone removed one blocker. ALOHAnet proved shared random radio
% access; the 1985 FCC ISM decision made spectrum permissionless; 802.11 (work began
% 1990, published 1997) replaced incompatible vendor islands; 1999 added 11a/b speed
% plus the WECA brand. First interoperability certifications followed in April 2000:
% https://www.wi-fi.org/news-events/newsroom/wireless-ethernet-compatibility-alliance-weca-awards-first-wi-fi
% Ask which milestone is regulatory, not technical: it is 1985, and it is arguably
% the most important one.
\begin{frame}{Four milestones made Wi-Fi possible}
  \vspace{0.3em}
  \centering
  % SOURCE: https://standards.ieee.org/wp-content/uploads/interactive/web/wi-fi-timeline/index.html
  % SOURCE (WECA formed 1999; first Wi-Fi certifications April 2000):
  % https://www.wi-fi.org/news-events/newsroom/wireless-ethernet-compatibility-alliance-weca-awards-first-wi-fi
  \begin{tikzpicture}[scale=0.9,transform shape]
    \draw[flow,very thick] (-0.6,0) -- (10.4,0);
    \foreach \x/\yr/\name/\what in {
      0/1971/ALOHAnet/{shared radio packet access},
      3.2/1985/{Unlicensed ISM}/{FCC opens spread-spectrum bands},
      6.4/1997/{IEEE 802.11}/{one MAC/PHY, no vendor islands},
      9.6/1999/{Wi-Fi}/{802.11a/b + the Wi-Fi brand (WECA)}}{
      \fill[ranblue] (\x,0) circle (2.2pt);
      \node[font=\bfseries\small,text=ranblue] at (\x,0.55) {\yr};
      \node[netnode,minimum width=2.2cm,minimum height=0.6cm,
            font=\bfseries\footnotesize] at (\x,-0.75) {\name};
      \node[font=\scriptsize,text=softgray,align=center,text width=2.5cm]
        at (\x,-1.75) {\what};
    }
  \end{tikzpicture}

  \vspace{0.8em}
  % The 1985 FCC decision opened the "junk bands" that microwave ovens pollute;
  % this hook pays off at the 802.11b three-clean-channels takeaway.
  \qa{Which of the four is regulatory rather than technical?}
     {1985. The FCC opened the ``junk bands'' microwave ovens pollute ---
      nobody needed permission to ship a radio there, and Wi-Fi built an
      empire on that permissionless ground.}

  \glossline{ISM = Industrial, Scientific, and Medical bands \quad
    FCC = Federal Communications Commission}
  \source{https://standards.ieee.org/wp-content/uploads/interactive/web/wi-fi-timeline/index.html}{IEEE 802.11 interactive timeline; Wi-Fi Alliance (WECA) newsroom, first certifications April 2000}
\end{frame}
